\documentclass[11pt]{article}\usepackage[margin=.85in]{geometry}\usepackage[T1]{fontenc}\begin{document}\section*{Independent UInt64 doubling review}
No blocking discrepancy was found at the stated source-word and normalized polynomial scope. The raw dbl expression preserves the source association and literal5.

\textbf{Boundary p=5.} The proof does not assume5 is canonical. It bounds $5a<2^{64}$ for canonical a and $p<2^{31}$, then proves raw multiplication modulo p equals multiplication by $5\bmod p$. The expression induction can therefore cover p=5 without changing the source body.

\textbf{Composition and guards.} All six outputs are canonical and decode to the field formula. The normalized theorem combines decoding with field validity under nonzero decoded z,w,y, nonzero2 and the point equation. Output nonzero is stated for decoded weights. Self-add dispatch excludes the earlier inverse guard; zero y selects inverse. No differently weighted converse is asserted.

\textbf{Independent checks.} Manifest hashes match. Fresh module, producer controls, separate axiom readback and reviewer controls exit0. Only propext, Classical.choice and Quot.sound occur. Additional controls check literal5 near the upper modulus bound, zero w with nonzero y selecting inverse, all nonzero coefficients in a canonical p7 payload, and the raw5 term at p5 versus p7. The full-premise p5 producer fixture is polynomial algebra, not smooth-curve evidence. Harmless linter warnings remain recorded.

Field.n=1, tables=None and ordinary execution/counter headroom remain external source conditions. No compiler/counter, divisor-class or whole-walk claim. No scientific trials or measured costs; producer unchanged. Exact hashes and receipts are in review.json and custody.json.
\end{document}